\section{Are Predictors that are Correlated with the Fundamental Process Selected More Often?}
\label{sec:fundamental_revisions_selection}

The theoretical extension in Section~\ref{sec:correlated_extension} predicts that predictors more strongly correlated with fundamentals should be selected more often. We now provide a direct empirical test of this prediction. The exercise asks whether the stock--predictor pairs for which a predictor is more closely related to innovations in firm fundamentals are also the stock--predictor pairs for which the rolling LASSO selects that predictor more frequently.

\paragraph{Data and identifying assumption.}

We proxy for innovations to firm fundamentals using analyst forecast revisions from I/B/E/S Detail History. The raw data are individual analyst EPS forecasts for U.S. firms between 2010 and 2017, matched to the CRSP PERMNOs in our LASSO estimation sample. We use quarterly EPS forecasts and keep the analyst-level panel before aggregation. A raw observation is indexed by firm $i$, analyst $a$, forecast announcement date $t$, and fiscal-quarter end $\tau$. For each firm--analyst--fiscal-quarter cell, we sort forecasts over time and define the revision
\[
\Delta EPS_{i,a,\tau,t}
= EPS_{i,a,\tau,t}-EPS_{i,a,\tau,t^-},
\]
where $t^-$ denotes the analyst's previous forecast for the same firm and fiscal quarter. We then scale the revision by the previous trading day's absolute CRSP price,
\[
\widetilde{\Delta EPS}_{i,a,\tau,t}
=\frac{\Delta EPS_{i,a,\tau,t}}{|P_{i,t-1}|}.
\]
Finally, we aggregate revisions to the firm-day level,
\[
FundNews_{i,t}
=\frac{1}{N_{i,t}}\sum_{a,\tau}\widetilde{\Delta EPS}_{i,a,\tau,t}.
\]

The identifying assumption is that analyst forecast revisions are timely, noisy signals of innovations to the market's information about firm fundamentals. Analysts revise EPS forecasts when new information changes their assessment of future cash flows. After conditioning on the analyst's own previous forecast for the same firm-quarter, the revision removes the analyst's prior level forecast and isolates the incremental news embedded in the update. Scaling by lagged price makes the revision comparable across firms. This proxy is imperfect: revisions may reflect analyst-specific processing, delayed incorporation of public news, or strategic forecast behavior. The identifying content of the test is therefore not that revisions are structural cash-flow shocks, but that they are monotone in the arrival of fundamental information relevant for firm value.

We date revisions using the I/B/E/S announcement date, mapped to the next CRSP trading day when the announcement falls on a weekend or holiday. This preserves the public-information timing interpretation while avoiding the loss of observations from non-trading-day announcements.

\paragraph{First stage: firm-specific fundamental relevance.}

Let $x_{j,t}$ denote predictor $j$ on trading day $t$. For each revision event date, we construct cumulative predictor innovations over an event window $W$ ending at or before the event date,
\[
X^W_{j,t}=\sum_{s\in W}x_{j,t+s}.
\]
We then estimate a firm-specific time-series regression separately for each firm $i$ and predictor $j$:
\[
FundNews_{i,t}
=a_i+\theta_{i,j}^{W}X^W_{j,t}+\varepsilon_{i,j,t},
\]
using only the revision-event dates of firm $i$. We require at least 50 revision-event observations for a firm to enter the baseline first stage. This threshold balances precision in the generated firm-specific relevance measure against cross-sectional coverage; below we report robustness to alternative thresholds. Fundamental relevance is defined as the absolute first-stage $t$-statistic,
\[
FundRel^W_{i,j}=|t(\hat\theta^W_{i,j})|,
\]
winsorized at the 1st and 99th percentiles within each specification.

For the baseline window $[-5,-1]$, the 50-event cutoff produces 160{,}740 firm--predictor relevance estimates, covering 846 firms and 190 predictors. In the broader event sample, the median firm contributes 165 usable revision-event days, with an interquartile range from 84 to 282.5. The cutoff therefore excludes firms with very sparse analyst-revision histories while retaining most firms with regular analyst coverage.

\paragraph{Second stage: selection rates and fundamental relevance.}

The second-stage outcome is the stock--predictor LASSO selection rate,
\[
SelectionRate_{i,j}
=\frac{1}{T_{i,j}}\sum_t \mathbf{1}\{\hat\beta_{i,j,t}\neq 0\},
\]
computed from the rolling LASSO coefficient tensor used in the main empirical analysis. We merge $SelectionRate_{i,j}$ with $FundRel^W_{i,j}$ and estimate
\[
SelectionRate_{i,j}
=\alpha_i+\delta_j+\beta FundRel^W_{i,j}+u_{i,j},
\]
where $\alpha_i$ are firm fixed effects and $\delta_j$ are predictor fixed effects. The coefficient $\beta$ is identified from within-firm, within-predictor variation in fundamental relevance. Standard errors are clustered two ways by firm and predictor.

Table~\ref{tab:fundamental_relevance_selection} reports the firm-specific results for four event windows. The mean selection rate in the estimation sample is 2.36 percent.

\begin{table}[H]
\centering
\caption{Fundamental Relevance and LASSO Selection}
\label{tab:fundamental_relevance_selection}
\begin{tabular}{lcccc}
\toprule
& \multicolumn{4}{c}{Dependent Variable: $SelectionRate_{i,j}$} \\
\cmidrule(lr){2-5}
Window & $[-5,-1]$ & $[-10,-1]$ & $[-5,0]$ & $[-10,0]$ \\
\midrule
$FundRel^W_{i,j}$
    & 0.000353** & 0.000333*** & 0.000395*** & 0.000343*** \\
    & (0.000144) & (0.000120) & (0.000132) & (0.000120) \\
\midrule
$t$-statistic
    & 2.45 & 2.77 & 2.99 & 2.86 \\
Standardized effect
    & 0.0056 & 0.0060 & 0.0064 & 0.0063 \\
Effect relative to mean selection
    & 1.52\% & 1.43\% & 1.70\% & 1.47\% \\
\midrule
Firm FE
    & Yes & Yes & Yes & Yes \\
Predictor FE
    & Yes & Yes & Yes & Yes \\
Observations
    & 160{,}740 & 160{,}740 & 160{,}740 & 160{,}740 \\
Firms
    & 846 & 846 & 846 & 846 \\
Predictors
    & 190 & 190 & 190 & 190 \\
\bottomrule
\end{tabular}

\vspace{0.15cm}
\begin{minipage}{0.95\textwidth}
\footnotesize
\textit{Note:} The table reports firm-specific event-window estimates requiring at least 50 usable revision-event observations per firm. The first stage regresses signed firm-day EPS forecast revision news, scaled by lagged price, on cumulative predictor innovations over the indicated window, separately for each firm and predictor. $FundRel^W_{i,j}$ is the absolute first-stage $t$-statistic, winsorized at the 1st and 99th percentiles. The second stage regresses the firm--predictor LASSO selection rate on $FundRel^W_{i,j}$ with firm and predictor fixed effects. Standard errors, in parentheses, are clustered two ways by firm and predictor. The standardized effect is $\hat\beta\,sd(FundRel^W)/sd(SelectionRate)$. The effect relative to mean selection is $\hat\beta/\overline{SelectionRate}$. $^{*}p<0.10$, $^{**}p<0.05$, $^{***}p<0.01$.
\end{minipage}
\end{table}

\paragraph{Robustness to the first-stage event-count cutoff.}

The baseline threshold of 50 revision-event days is meant to limit noise in the generated firm-specific relevance measure without restricting attention only to the most heavily covered firms. Table~\ref{tab:fundamental_relevance_cutoffs} reports the baseline $[-5,-1]$ specification under alternative minimum-event thresholds. The coefficient remains positive for every cutoff. Precision improves when moving from 20 to 50 or 100 events, but deteriorates once the cutoff becomes very high because the sample becomes much smaller and increasingly concentrated among large, heavily followed firms. We therefore use 50 events as the baseline and treat the other thresholds as robustness checks.

\begin{table}[H]
\centering
\caption{Robustness to Minimum Revision-Event Cutoff}
\label{tab:fundamental_relevance_cutoffs}
\resizebox{\textwidth}{!}{%
\begin{tabular}{lccccccc}
\toprule
Minimum events & 20 & 50 & 100 & 200 & 300 & 400 & 500 \\
\midrule
$FundRel^{[-5,-1]}_{i,j}$
    & 0.000310** & 0.000353** & 0.000394** & 0.000301 & 0.000544** & 0.000383 & 0.000223 \\
    & (0.000137) & (0.000144) & (0.000153) & (0.000186) & (0.000245) & (0.000338) & (0.000323) \\
$t$-statistic
    & 2.27 & 2.45 & 2.57 & 1.62 & 2.22 & 1.14 & 0.69 \\
\midrule
Firms
    & 924 & 846 & 676 & 391 & 208 & 118 & 78 \\
Observations
    & 175{,}560 & 160{,}740 & 128{,}440 & 74{,}290 & 39{,}520 & 22{,}420 & 14{,}820 \\
\bottomrule
\end{tabular}
}

\vspace{0.15cm}
\begin{minipage}{0.95\textwidth}
\footnotesize
\textit{Note:} The table repeats the baseline $[-5,-1]$ firm-specific specification while varying the minimum number of usable revision-event observations required for a firm to enter the first stage. Standard errors, in parentheses, are clustered two ways by firm and predictor. $^{*}p<0.10$, $^{**}p<0.05$, $^{***}p<0.01$.
\end{minipage}
\end{table}

\paragraph{Robustness to using coefficient magnitudes.}

Our baseline measure uses the absolute first-stage $t$-statistic because it captures the reliability of the firm-specific relation between predictor $j$ and fundamental news for firm $i$. The raw coefficient magnitude $|\hat\theta^W_{i,j}|$ is more directly interpretable as an exposure, but it also depends on the scale and volatility of firm-level EPS revisions. Thus, firms with noisier or more volatile forecast revisions may mechanically have larger coefficient magnitudes even when the predictor is not more reliably related to fundamentals. To check that the baseline result is not an artifact of using $|t|$, we re-estimate the second stage replacing $FundRel^W_{i,j}$ with the winsorized absolute coefficient $|\hat\theta^W_{i,j}|$ from the same firm-specific first-stage regressions.

\begin{table}[H]
\centering
\caption{Robustness: Using Absolute First-Stage Coefficients}
\label{tab:fundamental_relevance_selection_theta}
\begin{tabular}{lcccc}
\toprule
& \multicolumn{4}{c}{Dependent Variable: $SelectionRate_{i,j}$} \\
\cmidrule(lr){2-5}
Window & $[-5,-1]$ & $[-10,-1]$ & $[-5,0]$ & $[-10,0]$ \\
\midrule
$|\hat\theta^W_{i,j}|$
    & 0.4136* & -0.1069 & 0.0909 & 0.1298 \\
    & (0.2121) & (0.2567) & (0.1809) & (0.2333) \\
\midrule
$t$-statistic
    & 1.95 & -0.42 & 0.50 & 0.56 \\
Standardized effect
    & 0.0210 & -0.0039 & 0.0041 & 0.0046 \\
\midrule
Firm FE
    & Yes & Yes & Yes & Yes \\
Predictor FE
    & Yes & Yes & Yes & Yes \\
Observations
    & 160{,}740 & 160{,}740 & 160{,}740 & 160{,}740 \\
Firms
    & 846 & 846 & 846 & 846 \\
Predictors
    & 190 & 190 & 190 & 190 \\
\bottomrule
\end{tabular}

\vspace{0.15cm}
\begin{minipage}{0.95\textwidth}
\footnotesize
\textit{Note:} The table repeats the second-stage specification from Table~\ref{tab:fundamental_relevance_selection}, replacing $FundRel^W_{i,j}=|t(\hat\theta^W_{i,j})|$ with the winsorized absolute first-stage coefficient $|\hat\theta^W_{i,j}|$. Standard errors, in parentheses, are clustered two ways by firm and predictor. The standardized effect is $\hat\beta\,sd(|\hat\theta^W|)/sd(SelectionRate)$. $^{*}p<0.10$, $^{**}p<0.05$, $^{***}p<0.01$.
\end{minipage}
\end{table}

\paragraph{Discussion.}

The estimates are positive across all four windows. This is consistent with the mechanism in Section~\ref{sec:correlated_extension}: predictors that are more closely related to innovations in firm fundamentals are selected more frequently by the rolling LASSO. The baseline pre-event window $[-5,-1]$ yields a coefficient of 0.00035 with a two-way clustered $t$-statistic of 2.45. Extending the pre-event window to $[-10,-1]$ gives a similar estimate, 0.00033, with a $t$-statistic of 2.77. Including the event day strengthens the coefficient for the five-day window: the estimate rises to 0.00040 for $[-5,0]$ and the $t$-statistic increases to 2.99. The corresponding ten-day window $[-10,0]$ remains positive and statistically significant.

The economic magnitudes are small in levels, which is expected because unconditional LASSO selection rates are low. A one-unit increase in fundamental relevance raises the selection rate by roughly 1.4--1.7 percent of the mean selection rate, depending on the window. Standardized effects are also modest, around 0.6 percent of a standard deviation in selection rates. The results should therefore not be read as implying that analyst-revision fundamentals are the dominant determinant of selection. Rather, they show that the identity of selected predictors is not purely arbitrary: within firm and within predictor, the pairs with stronger links to measured fundamental-news revisions are selected more often.

The timing pattern is informative. The effect is present using only pre-event predictor innovations, which mitigates the concern that the relation is entirely contemporaneous. At the same time, the stronger $[-5,0]$ estimate suggests that part of the relevant information arrives on the revision day itself, consistent with analysts updating forecasts in response to same-day public news. Table~\ref{tab:fundamental_relevance_cutoffs} shows that the baseline result is not driven by the 50-event threshold: the coefficient remains positive for cutoffs from 20 to 500 events, although precision falls at very high cutoffs as the sample shrinks. The coefficient-based robustness check in Table~\ref{tab:fundamental_relevance_selection_theta} is weaker: the baseline $[-5,-1]$ estimate remains positive and marginally significant, but the longer-window coefficient specifications are not stable. This pattern is consistent with the interpretation that raw coefficient magnitudes are contaminated by cross-firm differences in the volatility of EPS revisions, whereas the $t$-statistic better captures the reliability of the predictor--fundamental-news relation. Overall, the evidence supports the paper's interpretation that finite-sample LASSO selection is amplified when candidate predictors are genuinely correlated with the firm's fundamental information process, while also suggesting that the signal is measured with substantial noise.
